\documentclass[fontsize=15pt]{jlreq}

%プリアンブル
\usepackage{amsmath}
\usepackage{mathtools}
\pagestyle{empty}
\usepackage[top=15mm, bottom=20mm, left=20mm, right=20mm]{geometry}

%コマンド
\newcommand{\m}[1]{\raisebox{0.1em}{\textcircled{\scriptsize #1}}}



%本文
\begin{document}
\begin{center}
{\LARGE {振り返り}} \\
-生殖- %単元ごとに変えるか消す
\end{center}
\vspace{1em}
\begin{flushright}
名前：\underline{\hspace{5cm}}
\end{flushright}

% 問題部分
\begin{enumerate}

    \item 有性生殖と無性生殖をそれぞれ説明せよ．
    \begin{itemize}
        \item 有性生殖：
        \vspace{1em}
        \item 無性生殖：
    \end{itemize}
    \vspace{2em} 

    \item 植物は有性生殖だけでなく，無性生殖も行う．この無性生殖を特になんというか．\\
    (\hspace{3cm})
    \vspace{2em}

    \item 次の文を穴埋めせよ．
    \begin{itemize}
        \item 動物の受精：雌が作る(\hspace{2cm})と雄が作る(\hspace{2cm})が受精して，受精卵ができる．
        \vspace{1em}
        \item 被子植物の受精：花粉がめしべの柱頭につくと，(\hspace{3cm})が胚珠へと伸びていく．花粉管の中の(\hspace{3cm})と胚珠の中の(\hspace{3cm})が受精して受精卵ができる．
    \end{itemize}
    \vspace{2em} 

    \item 受精卵になってから個体の組織などが形成されるまでの状態をなんというか．\\
    (\hspace{2cm})
    \vspace{2em} 

    \item 受精してから個体として体のつくりが発達していく過程をなんというか．\\
    (\hspace{2cm})
    \vspace{2em}   

\end{enumerate}

\end{document}
